\section{为什么数学层必须扩展，但不能漂移}

本手稿的早期版本使用势阱图景，是因为它便于教学。
对于第一轮写作来说，这是合适的；但现在已经不够。
一旦本书开始严肃讨论意识状态转变、框架形成与干预设计，
数学层就必须更明确地说明：我们究竟在设想什么类型的系统。

与此同时，本章必须继续服从证据纪律。
这里的大部分内容都属于\textbf{中等}车道。
本章借助文献地图中的高级动力系统来源，提升解释精度：
包括 Nature Reviews Neuroscience 关于大脑中吸引子与积分器网络的综述、
上一章已经使用过的亚稳态综述，以及本地索引的关于自组织吸引子网络与 Koopman 式视角的 arXiv 论文。
这些来源帮助手稿更清楚地讨论持续、切换、振荡与反馈。
它们\emph{并不}证明人的日常常规在字面上就是一个已知参数的低维微分方程。

\begin{discussion}
数学深度的业务价值因此是狭义而具体的。
它应当回答的问题包括：
\begin{enumerate}[nosep]
  \item 什么使一个状态在局部上稳定；
  \item 什么使一个转变是突变式而非渐进式；
  \item 为什么某些模式会以节律方式重复出现；
  \item 为什么历史的反复积累会使重新进入更容易；
  \item 干预为何可以被解释为控制输入，而不是道德劝诫。
\end{enumerate}
如果某个数学工具不能锐化上述问题之一，它就不应出现在本章正文之中。
\end{discussion}